\documentclass[english,10pt,table]{beamer}
%\documentclass[english,10pt,handout]{beamer}
%\documentclass[slidestop,usepdftitle=false]{beamer}

% Copyright 2007 by Till Tantau
%
% This file may be distributed and/or modified
%
% 1. under the LaTeX Project Public License and/or
% 2. under the GNU Public License.
%
% See the file doc/licenses/LICENSE for more details.


% Common packages

\usepackage[english]{babel}
\usepackage[utf8]{vietnam}
%\usepackage{times}
\usefonttheme[onlymath]{serif}
\usecolortheme{default}
\usepackage{booktabs}
\usepackage{mathpartir}
\usepackage{listings}
\usepackage{pbox}
\mprset{flushleft}
\mode<article>
{
  \usepackage{times}
  \usepackage{mathptmx}
  \usepackage[left=1.5cm,right=6cm,top=1.5cm,bottom=3cm]{geometry}
}

\usepackage{hyperref}


\usepackage{tikz}
\usetikzlibrary{arrows,snakes,backgrounds}
\tikzstyle{mnode}=[circle, draw, fill=black, inner sep=0pt, minimum width=4pt]
\usepackage{colortbl}
%\usepackage{yfonts}
\usepackage{translator} % comment this, if not available


% Common settings for all lectures in this course

\def\lecturename{Khai phá dữ liệu}

\title{\insertlecture}

\author{Nguyễn Giáp Nguyên Sinh \\
		Nguyễn Quốc Long}

\institute
{
  Khoa Khoa học và Kỹ thuật Máy tính\\
  Đại học Bách Khoa TP.HCM
}

\subject{Lecturer \lecturename}




% Beamer version theme settings

\useoutertheme[height=0pt,width=2cm,right]{sidebar}
\usecolortheme{rose,sidebartab}
\useinnertheme{circles}
\usefonttheme[only large]{structurebold}

\setbeamercolor{sidebar right}{bg=black!15}
\setbeamercolor{structure}{fg=blue}
\setbeamercolor{author}{parent=structure}

\setbeamerfont{title}{series=\normalfont,size=\LARGE}
\setbeamerfont{title in sidebar}{series=\bfseries}
\setbeamerfont{author in sidebar}{series=\bfseries}
\setbeamerfont*{item}{series=}
\setbeamerfont{frametitle}{size=}
\setbeamerfont{block title}{size=\small}
\setbeamerfont{subtitle}{size=\normalsize,series=\normalfont}

\setbeamertemplate{navigation symbols}{}
\setbeamertemplate{bibliography item}[book]
\setbeamertemplate{sidebar right}
{
  {\usebeamerfont{title in sidebar}%
    \vskip1.5em%
    \hskip3pt%
    \usebeamercolor[fg]{title in sidebar}%
    %\insertshorttitle[width=2cm,center,respectlinebreaks]\par%
    \vskip1.25em%
  }%
  {%
    \hskip3pt%
    \usebeamercolor[fg]{author in sidebar}%
    \usebeamerfont{author in sidebar}%
    %\insertshortauthor[width=2cm,center,respectlinebreaks]\par%
    \vskip1.25em%
  }%
  \hbox to2cm{\hss\insertlogo\hss}
  \vskip1.25em%
  \insertverticalnavigation{2cm}%
  \vfill
  \hbox to 2cm{\hfill\usebeamerfont{subsection in
      sidebar}\strut\usebeamercolor[fg]{subsection in
      sidebar}Ch.\insertshortlecture, p.\insertframenumber/\inserttotalframenumber\hskip5pt}%
  \vskip3pt%
}%

\setbeamertemplate{title page}
{
  \vbox{}
  \vskip1em
  {\huge Chapter \insertshortlecture\par}
  {\usebeamercolor[fg]{title}\usebeamerfont{title}\inserttitle\par}%
  \ifx\insertsubtitle\@empty%
  \else%
    \vskip0.25em%
    {\usebeamerfont{subtitle}\usebeamercolor[fg]{subtitle}\insertsubtitle\par}%
  \fi%     
  \vskip1em\par
  %\emph{\lecturename}\ on \insertdate\par
  \vskip0pt plus1filll
  \leftskip=0pt plus1fill\insertauthor\par
  \insertinstitute\vskip1em
}

\logo{\includegraphics[width=1.5cm]{hcmut.png}}



% Article version layout settings

\mode<article>

\makeatletter
\def\@listI{\leftmargin\leftmargini
  \parsep 0pt
  \topsep 5\p@   \@plus3\p@ \@minus5\p@
  \itemsep0pt}
\let\@listi=\@listI


\setbeamertemplate{frametitle}{\paragraph*{\insertframetitle\
    \ \small\insertframesubtitle}\ \par
}
\setbeamertemplate{frame end}{%
  \marginpar{\scriptsize\hbox to 1cm{\sffamily%
      \hfill\strut\insertshortlecture.\insertframenumber}\hrule height .2pt}}
\setlength{\marginparwidth}{1cm}
\setlength{\marginparsep}{4.5cm}

\def\@maketitle{\makechapter}

\def\makechapter{
  \newpage
  \null
  \vskip 2em%
  {%
    \parindent=0pt
    \raggedright
    \sffamily
    \vskip8pt
    {\fontsize{36pt}{36pt}\selectfont Chapter \insertshortlecture \par\vskip2pt}
    {\fontsize{24pt}{28pt}\selectfont \color{blue!50!black} \insertlecture\par\vskip4pt}
    {\Large\selectfont \color{blue!50!black} \insertsubtitle\par}
    \vskip10pt

    \normalsize\selectfont Print version of
    Lecturer \emph{\lecturename} of \@date\par\vskip1.5em
    \hfill Tran Vinh Tan, Faculty of CSE, University of Technology
  }
  \par
  \vskip 1.5em%
}

\let\origstartsection=\@startsection
\def\@startsection#1#2#3#4#5#6{%
  \origstartsection{#1}{#2}{#3}{#4}{#5}{#6\normalfont\sffamily\color{blue!50!black}\selectfont}}

\makeatother

\mode
<all>




% Typesetting Listings

\usepackage{listings}
\lstset{language=Java}

\alt<presentation>
{\lstset{%
  basicstyle=\footnotesize\ttfamily,
  commentstyle=\slshape\color{green!50!black},
  keywordstyle=\bfseries\color{blue!50!black},
  identifierstyle=\color{blue},
  stringstyle=\color{orange},
  escapechar=\#,
  emphstyle=\color{red}}
}
{
	\lstset{%
		basicstyle=\ttfamily,
		keywordstyle=\bfseries,
		commentstyle=\itshape,
		escapechar=\#,
		emphstyle=\bfseries\color{red}
	}
}



% Common theorem-like environments

\theoremstyle{example}
\newtheorem{exercise}[theorem]{\translate{Exercise}}


% New useful definitions:

\newbox\mytempbox
\newdimen\mytempdimen

\newcommand\includegraphicscopyright[3][]{%
	\leavevmode\vbox{\vskip3pt\raggedright\setbox\mytempbox=\hbox{\includegraphics[#1]{#2}}%
		\mytempdimen=\wd\mytempbox\box\mytempbox\par\vskip1pt%
		\fontsize{3}{3.5}\selectfont{\color{black!25}{\vbox{\hsize=\mytempdimen#3}}}\vskip3pt%
}}

\newenvironment{colortabular}[1]{\medskip\rowcolors[]{1}{blue!20}{blue!10}\tabular{#1}\rowcolor{blue!40}}{\endtabular\medskip}

\def\equad{\leavevmode\hbox{}\quad}

\newenvironment{greencolortabular}[1]
{\medskip\rowcolors[]{1}{green!50!black!20}{green!50!black!10}%
	\tabular{#1}\rowcolor{green!50!black!40}}%
{\endtabular\medskip}




\lecture[1]{Elementary analysis}{lecture-text}

%\date[]{~~}
%%%%%%%%%%%%%%%%%%%%%%%%%%%%%%%%%%%%%%%%%%%%%%%%%%%%%%%%%%%%%%%%%%%%%
%%%%%%%%%%%%%%%%%%%%%%%%%%%%%%%%%%%%%%%%%%%%%%%%%%%%%%%%%%%%%%%%%%%%%
\begin{document}

\begin{frame}
  \maketitle
\end{frame}

\section*{Contents}
\begin{frame}\frametitle<presentation>{Contents}
  \tableofcontents
\end{frame}


%%%%%%%%%%%%%%%%%%%%%%%%%%%
\section{Definition and notations}
\subsection{Algorithm}
\begin{frame}{Algorithm}		
	\begin{block}{What is an algorithm?}
	 An algorithm is a finite set of precise instructions for performing a computation or for solving a problem.
  \end{block}	
\pause
		
	\begin{block}{Properties of algorithms}
	 \begin{itemize}
		 \item \alert{Input} from a specified set,
		 \item \alert{Output} from a specified set (solution),
		 \item \alert{Definiteness} of every step in the computation,
		 \item \alert{Correctness} of output for every possible input,
		 \item \alert{Finiteness} of the number of calculation steps,
		 \item \alert{Effectiveness} of each calculation step and
		 \item \alert{Generality} for a class of problems.
   \end{itemize}
  \end{block}	
\end{frame}

%%%%%%%%%%%%%%%%%%%%%%%%%%%%
\subsection{Complexity}
\begin{frame}{Complexity}
		
	\begin{block}{}
	 	 \begin{itemize}\small
		 \item Generally, not much interested in time and space complexity for small inputs.
		 \item Given two algorithms \alert{$A$} and \alert{$B$} for solving problem $P$.
		\end{itemize}
  \end{block}	
	
	\begin{center}
	\small\begin{tabular}{|c|c|c|}
			\hline
			Input size & Algorithm A & Algorithm B \\ \hline
		$n$ & \textcolor{blue}{5000 $n$} & \textcolor{red}{$1.2^n$}\\ \hline
		$10$ & \textcolor{blue}{50,000} & \textcolor{red}{6}\\ \hline
		$100$ & \textcolor{blue}{500,000} & \textcolor{red}{2,817,975}\\ \hline
		$1,000$ & \textcolor{blue}{5,000,000} & \textcolor{red}{$1.5 \times 10^{79}$}\\ \hline
		$100,000$ & \textcolor{blue}{$5 \times 10^8$} & \textcolor{red}{$1.3 \times 10^{7918}$}\\ \hline
	\end{tabular}
 \end{center}
	
	\begin{block}{}
	 	 \begin{itemize}\small
		 \item <2-> \alert{$B$} cannot be used for large inputs, while \alert{$A$} is still feasible.
		 \item <2-> So what is important is the \alert{growth} of the complexity functions.
		 \item <2-> Growth of time and space complexity with increasing input size \alert{$n$} is a suitable measure for the comparison of algorithms. 
		\end{itemize}
  \end{block}		
\end{frame}

%%%%%%%%%%%%%%%%%%%%%%%%%%%%
\subsection{Formulas}
\begin{frame}{Types of formulas for basic operation count}
		
	\begin{block}{}
	 	 \begin{itemize}\small
		 \item Exact formulas, e.g., \alert{$C(n) = n(n-1)/2$}.
		 \item <2-> Formula indicating order of growth with specific multiplicative constant
            e.g., \alert{$C(n) \approx 0.5 n^2$}.
		 \item <3-> Formula indicating order of growth with unknown multiplicative constant
            e.g., \alert{$C(n) \approx c.n^2$}
		 \item <4-> Most important: Order of growth within a constant multiple as \alert{$n \rightarrow \infty$}
		\end{itemize}
  \end{block}	
 \only<5>{	
	\begin{block}{Asymptotic growth rate}
	 	 A way of comparing functions that ignores constant factors and small input sizes
		\begin{itemize}\small
		 \item \alert{$O(g(n))$}: class of functions \alert{$f(n)$} that grow \alert{no faster than} \alert{$g(n)$}
 		 \item \alert{$\Theta(g(n))$}: class of functions \alert{$f(n)$} that grow \alert{at the same rate as} \alert{$g(n)$}
		 \item \alert{$\Omega(g(n))$}: class of functions \alert{$f(n)$} that grow \alert{at least as fast as} \alert{$g(n)$}
		\end{itemize}
  \end{block}	
 }
\end{frame}


%%%%%%%%%%%%%%%%%%%%%%%%%%%%
\begin{frame}{Complexity classes - a small vocabulary}
	\begin{block}{}
	 \begin{itemize}\small
 		  \item \textcolor{blue}{Constant}: \alert{$O(1)$} (independing on the input size)
     \item \textcolor{blue}{Sub-linear} or \textcolor{blue}{logarithmic}: \alert{$O(\log n)$}
     \item \textcolor{blue}{Linear}: \alert{$O(n)$}
					\item \textcolor{blue}{Quasi-linear}: \alert{$O(n \log n)$}
					\item \textcolor{blue}{Quadratic}: \alert{$O(n^2)$}
					\item \textcolor{blue}{Cubic}: \alert{$O(n^3)$}
     \item \textcolor{blue}{Polynomial}: \alert{$O(n^p)$} (\alert{$O(n^2)$}, \alert{$O(n^3)$}, etc)
     \item \textcolor{blue}{Quasi-polynomial}: \alert{$O(n^{\log(n)})$}
     \item \textcolor{blue}{Exponential}: \alert{$O(2^n)$}
					\item \textcolor{blue}{Factorial}: \alert{$O(n!)$}
		\end{itemize}
 \end{block}	
\end{frame}


\begin{frame}{Asymptotic upper bound - worst case}
	
	\begin{block}{Asymptotic upper bound  ``big O''}\small
		\textcolor{blue}{$T(n)$ = \alert{$O$}$(f(n))$} \textbf{iif} 
		\alert{$\exists c \in R^+$}, \alert{$c>0$} and \alert{$\exists n_0 \in N$}, \alert{$n_0>0$} \\
		such that \alert{$\forall n>n_0$}: \alert{$T(n) \leq c \times f(n) $}
 \end{block}	
	\pause
		\begin{block}{Example}\small
		Let \alert{$T(n) = 2n+3n^3 + 5$}. \alert{$T(n)$} is in \alert{$O(n^3)$} with:
	 \begin{itemize}
     \item (\alert{$c=8$} and \alert{$n_0=1$}) or (\alert{$c=5$} and \alert{$n_0=2$})
		\end{itemize}
 \end{block}	
	\pause
	\begin{block}{}\small
		\textbf{Principle}: \textcolor{blue}{the lower-order terms are negligible.}
 \end{block}		
\end{frame}


\begin{frame}{Asymtotic lower bound - best case}
	
	\begin{block}{``big Omega''}\small
		\textcolor{blue}{$T(n)$ = \alert{$\Omega$}$(f(n))$} \textbf{iif} 
		\alert{$\exists c \in R^+$}, \alert{$c>0$} and \alert{$\exists n_0 \in N$}, \alert{$n_0>0$}\\
		such that \alert{$\forall n>n_0$}: \alert{$T(n) \geq c \times f(n) $}
 \end{block}	
	\pause
		\begin{block}{Example}\small
		Let \alert{$T(n) = 2n+3n^3 + 5$}. \alert{$T(n)$} is in \alert{$\Omega(n^3)$} with:
	 \begin{itemize}
     \item (\alert{$c=1$} and \alert{$n_0=1$}) 
		\end{itemize}
 \end{block}	
\end{frame}


\begin{frame}{Asymtotic approximating bound - average case}
	
	\begin{block}{``big Theta''}
  \textcolor{blue}{$T(n)$ = \alert{$\Theta$}$(f(n))$} \textbf{iif}
		\alert{$\exists c_1, c_2 \in R^+$}, \alert{$c_1>0$}, \alert{$c_2>0$} and \alert{$\exists n_0 \in N$}, \alert{$n_0>0$}\\
		such that \alert{$\forall n>n_0$}: \alert{$c_1 \times f(n) \leq T(n) \leq c_2 \times f(n) $}
 \end{block}	
		\pause
		\begin{block}{Property}\small
		\alert{$T(n) = O(f(n))$} and \alert{$T(n) = \Omega(f(n))$} $\Longrightarrow$ \alert{$T(n) = \Theta(f(n))$}
 \end{block}	
	\pause
	\begin{block}{Example}\small
		Let \alert{$T(n) = 2n+3n^3 + 5$}. \\
		So, \alert{$T(n)$} is in \alert{$O(n^3)$} and in \alert{$\Omega(n^3)$}.\\
		Consequently, \alert{$T(n)$} is in \alert{$\Theta(n^3)$}.
 \end{block}	
\end{frame}

%%%%%%%%%%%%%%%%%%%%%%%%%%%%
\begin{frame}{Other properties}
	
	\begin{block}{Not transitive}\footnotesize
	 \begin{itemize}
     \item $f(n) = n^2$; $g(n)= n$
		 \item $\Rightarrow f(n) = O(n^2) = g(n)$ but $f(n) \neq g(n)$
	\end{itemize}
 \end{block}
\pause
   \begin{block}{Transitivity}\small				
			\begin{itemize}\small				
		   \item $f(n) = O(g(n))$ \& $g(n) = O(h(n))$ $\Rightarrow f(n) = O(h(n))$
			 \item $f(n) = \Omega(g(n))$ \& $g(n) = \Omega(h(n))$ $\Rightarrow f(n) = \Omega(h(n))$
			 \item $f(n) = \Theta(g(n))$ \& $g(n) = \Theta(h(n))$ $\Rightarrow f(n) = \Theta(h(n))$
      \end{itemize}
   \end{block}
\pause
    \begin{block}{Additivity}\small				
			\begin{itemize}\small				
		   \item $f(n) = O(h(n))$ \& $g(n) = O(h(n))$ $\Rightarrow f(n) + g(n) = O(h(n))$
			 \item $f(n) = \Omega(h(n))$ \& $g(n) = \Omega(h(n))$ $\Rightarrow f(n) + g(n) = \Omega(h(n))$
			 \item $f(n) = \Theta(h(n))$ \& $g(n) = \Theta(h(n))$ $\Rightarrow f(n) + g(n) = \Theta(h(n))$
      \end{itemize}
   \end{block}
\end{frame}

%%%%%%%%%%%%%%%%%%%%%%%%%%%%
\begin{frame}{Exercise}
	
	\begin{block}{}\footnotesize
	 \begin{itemize}
     \item <2->\alert{$T(n) = 3 + 5n^2$} $\Rightarrow$ \alert{$T(n) = \Theta(n^2)$} ?
		 \item <3->if \alert{$T(n) = 
		    \left\{ \begin{array}{cccc} 
			    2n + 5 &  if & n & \mbox{ is even }\\
				  n^2 - n + 1 &  if & n & \mbox{ is odd }\\
								\end{array} \right.$}, then \alert{$T(n) = O(?)$} and \alert{$T(n) = \Omega(?)$}.
		\end{itemize}
 \end{block}	
\end{frame}
%%%%%%%%%%%%%%%%%%%%%%%%%%%%
\begin{frame}{Exercise}
	
	\begin{block}{}
  Compare the asymptotic behaviours of 
	\begin{enumerate}
		\item \alert{$2^n$} and \alert{$10^n$}
		\item \alert{$\log_{2}n$} and \alert{$\log_{10}n$}
	\end{enumerate}	
 \end{block}

\pause

	\begin{block}{}
	\begin{enumerate}
  \item Prove that for any positive functions \alert{$f$} and \alert{$g$}, \alert{$f(n) +g(n)$} and \alert{$max(f(n); g(n))$} are asymptotically equivalent. \\
	\item Give a (necessary and sufficient) condition on positive functions \alert{$f$} and \alert{$g$} to ensure that \alert{$f(n) +g(n)$} and \alert{$f(n)$} are asymptotically equivalent. 
	\end{enumerate}
	\end{block}	
\end{frame}

%%%%%%%%%%%%%%%%%%%%%%%%%%%%

\begin{frame}{Common asymptotic behaviours}
	\begin{block}{}\footnotesize
	\begin{tabular}{|c|c|c|c|c|c|c|}
		\hline
		Size & \multicolumn{6}{c|}{Approximate computational time}\\ \hline
		$n$ & $\Theta(\log n)$ & $\Theta(n)$ & $\Theta(n \log n)$ & $\Theta(n^2)$ & $\Theta(2^n)$ & $\Theta(n!)$ \\ \hline
		$10$ & $3. 10^{-9}$s & $10^{-8}$s & $3. 10^{-8}$s & $10^{-7}$s & $10^{-6}$s & $3. 10^{-3}$s\\
		$10^2$ & $7. 10^{-9}$s & $10^{-7}$s & $7. 10^{-7}$s & $10^{-5}$s & $4. 10^{13}$y  & * \\
		$10^3$ & $10^{-8}$s & $10^{-6}$s & $10^{-5}$s & $10^{-3}$s  & * & * \\
		$10^4$ & $1,3. 10^{-8}$s & $10^{-5}$s & $10^{-4}$s & $10^{-1}$s & * & * \\
		$10^5$ & $1,7. 10^{-8}$s & $10^{-4}$s & $2. 10^{-3}$s & $10$s & * & * \\
		$10^6$ & $2. 10^{-8}$s & $10^{-3}$s & $2. 10^{-2}$s & $17$m & * & * \\ \hline
	\end{tabular}
  \end{block}
\end{frame}



%%%%%%%%%%%%%%%%%%%%%%%%%%%%
\section{Basic methods for asymptotic behaviour analysis}
\subsection{Counting number of elementary operations}
\begin{frame}{Example}
		
	\begin{block}{}
	 \begin{enumerate}[1.]\small
		   \item \only<2->{\textcolor{orange}{(1)}}\textbf{Var} int: \alert{$d=0$} 
	    \item \only<3->{\textcolor{orange}{($n$)}}\textbf{For} \alert{$i$} from \alert{$1$} to \alert{$n$} \textbf{do}
				 	\begin{enumerate}[1.]\small
	      \item \only<4->{\textcolor{orange}{(1)}}\alert{$d$} = \alert{$d + 1$}
	      \item \only<5->{\textcolor{orange}{(1)}}\alert{$a[i]$} = \alert{$a[i] \times a[i] + d \times d$}
						\end{enumerate}
					\item \textbf{Endfor}
    \end{enumerate}
  \end{block}	
		\only<6->{
		\begin{block}{}
		 Number of elementary operations: \alert{$1 + n \times ( 1 + 1) = 2n +1$}.
	 \end{block}	
		}
\end{frame}
%%%%%%%%%%%%%%%%%%%%%%%%%%%
\begin{frame}{Linear loop example}
  \begin{columns}
  \begin{column}{0.48\textwidth}
   \begin{block}{}\small				
			\begin{enumerate}[1.]\small				
		   \item \textbf{Var} int: \alert{$i=1$}
     \item \textbf{While} \alert{$i \leq n$} \textbf{do}
					 \begin{enumerate}[1.]\small				
       \item Write ``Bonjour'' 
							\item \alert{$i = i + 1$}
						\end{enumerate}
					\item \textbf{EndWhile}
			\end{enumerate}
   \end{block}
	  \end{column}
			\pause
   \begin{column}{0.48\textwidth}
    \begin{block}{}\small				
			\begin{enumerate}[1.]\small				
		   \item \textbf{Var} int: \alert{$i=n$}
     \item \textbf{While} \alert{$i \geq 1$} \textbf{do}
					 \begin{enumerate}[1.]\small				
       \item Write ``Bonjour'' 
							\item \alert{$i = i - 1$}
						\end{enumerate}
					\item \textbf{EndWhile}
			\end{enumerate}
   \end{block}
  \end{column}
  \end{columns}
	
		\only<3->{
		\begin{block}{}
		 Number of elementary operations: \alert{$2n+1$}.
	 \end{block}	
		}
\end{frame}

%%%%%%%%%%%%%%%%%%%%%%%%%%%
\begin{frame}{Logarithmic loop example}
  \begin{columns}
  \begin{column}{0.48\textwidth}
   \begin{block}{}\small				
			\begin{enumerate}[1.]\small				
		   \item \textbf{Var} int: \alert{$i=1$}
     \item \textbf{While} \alert{$i \leq n$} \textbf{do}
					 \begin{enumerate}[1.]\small				
       \item Write ``Bonjour'' 
							\item \alert{$i = i \times 2$}
						\end{enumerate}
					\item \textbf{EndWhile}
			\end{enumerate}
   \end{block}
	  \end{column}
			\pause
   \begin{column}{0.48\textwidth}
    \begin{block}{}\small				
			\begin{enumerate}[1.]\small				
		   \item \textbf{Var} int: \alert{$i=n$}
     \item \textbf{While} \alert{$i \geq 1$} \textbf{do}
					 \begin{enumerate}[1.]\small				
       \item Write ``Bonjour'' 
							\item \alert{$i = i / 2$}
						\end{enumerate}
					\item \textbf{EndWhile}
			\end{enumerate}
   \end{block}
  \end{column}
  \end{columns}
		
		\only<3->{
		\begin{block}{}
		 Number of elementary operations: \alert{$1 + \log_2 (n)$}.
	 \end{block}	
		}
\end{frame}

%%%%%%%%%%%%%%%%%%%%%%%%%%%
\begin{frame}{Nested loop example}
   \only<1->{
			\begin{block}{}\small			
     \textcolor{blue}{Nb of iterations = nb of iterations of external loop $\times$ nb of iterations of internal loop}
			\end{block}
			}
			\pause
			\begin{columns}
  \begin{column}{0.48\textwidth}			
    \begin{block}{}\small				
			\begin{enumerate}[1.]\small				
		   \item \textbf{Var} int: \alert{$i=1$}
     \item \textbf{While} \alert{$i \leq n$} \textbf{do}
					 \begin{enumerate}[1.]\small				
       \item \textbf{Var} int: \alert{$j=1$}
							\item \textbf{While} \alert{$j \leq n$} \textbf{do}
					   \begin{enumerate}[1.]\small				
         \item Write ``Bonjour'' 
									\item \alert{$j = j \times 3$}
							 \end{enumerate}
					  \item \textbf{EndWhile}
							\item \alert{$i = i + 1$}
						\end{enumerate}
					\item \textbf{EndWhile}
			\end{enumerate}
   \end{block}
		 \end{column}
			\pause
   \begin{column}{0.48\textwidth}
		\begin{block}{}
		 Number of elementary operations: \alert{$1+ n + n \times \log_3 (n)$}.
	 \end{block}	
  \end{column}
  \end{columns}
\end{frame}

%%%%%%%%%%%%%%%%%%%%%%%%%%%
\begin{frame}{Exercise}
	\begin{block}{Function XYZ(array: $a[]$)}\small				
			\begin{enumerate}[1.]\small				
		   \item \only<2->{\textcolor{orange}{(1)}}\textbf{Var} int: \alert{$i$}
     \item \only<3->{\textcolor{orange}{($n$)}}\textbf{For} \alert{$i$} from \alert{$1$} to \alert{$n$} \textbf{do}
					 \begin{enumerate}[1.]\small				
       \item \only<4->{\textcolor{orange}{(1)}}\textbf{Var} int: \alert{$t$} = \alert{$a[i]$}
							\item \only<5->{\textcolor{orange}{(1)}}\textbf{Var} int: \alert{$j$}
							\item \only<6->{\textcolor{orange}{(?)}}\textbf{For} \alert{$j$} from \alert{$i-1$} to \alert{$0$} \textbf{do}
							 \begin{enumerate}[1.]\small				
         \item \only<7->{\textcolor{orange}{(1)}}\alert{$a[j+1]$} = \alert{$a[j]$}
								\end{enumerate}
					  \item \textbf{EndFor}
							\item \only<8->{\textcolor{orange}{(1)}}\alert{$a[j+1]$} = \alert{$t$}
						\end{enumerate}
					\item \textbf{EndFor}
			\end{enumerate}
 \end{block}		
\end{frame}
%%%%%%%%%%%%%%%%%%%%%%%%%%%%



%%%%%%%%%%%%%%%%%%%%%%%%%%%%
\begin{frame}{Homeworks}
	
	\begin{block}{}\footnotesize
	\textcolor{blue}{Give algorithms having number of elementary operations as below.}\\
	 \begin{itemize}
     \item \alert{$T_1(n) = 3 + 5n$}
					\item \alert{$T_2(n) = n \log_2 n$}
					\item \alert{$T_3(n) = n^3$}
					\item \alert{$T_4(n) = (3n)!$}
					\item \alert{$T_5(n) =  \log_2 (3 n) $}
					\item \alert{$T_6(n) = 2 \log_3 (2 n)$}
					\item \alert{$T_7(n) = n^2 \log_4 n $}
					\item \alert{$T_8(n) = \sqrt{n} $}
					\item \alert{$T_{9}(n) = \sqrt[3]{n^2} $}
					\item \alert{$T_{10}(n) = 2^n $}
					\item \alert{$T_{11}(n) = n! $}
		\end{itemize}
 \end{block}	
\end{frame}
%%%%%%%%%%%%%%%%%%%%%%%%%%%%

\end{document}